\section*{Limitations}

While our results show that a self-speculating agent can improve next-tool-call
prediction while preserving downstream task performance, the present study is an
initial step rather than a complete treatment of speculative tool use. Our
method and evaluation still have several limitations that should be considered
when interpreting the results.

Our method assumes that speculative tool calls can be issued without changing the
external environment. This is appropriate for read-only tools such as search,
retrieval, database lookup, and many information-gathering API calls, where an
incorrect speculation can simply be discarded. It is less suitable for tools
that mutate external state, such as placing an order, updating a database
record, sending a message, or triggering an irreversible workflow. In such
settings, speculative execution would require additional safeguards, such as
dry-run modes, transactional rollback, human confirmation, or restricting
speculation to the read-only prefix of a task. Since a large fraction of agentic
tool use is read-only or has a read-only planning stage, we expect speculation
to remain broadly useful, but the deployment policy must distinguish safe calls
from state-changing calls.

Our empirical evaluation is also limited in scope. We evaluate self-speculation
on two agentic task families: multi-hop search question answering and
conversational tool use. These settings cover different tool-call patterns,
but they do not exhaust the diversity of agent environments, especially code
execution, web navigation, long-running workflows, and multi-agent tool use.
Moreover, our training experiments are conducted with 4B-scale backbone models.
Larger models may exhibit different optimization dynamics, different
speculator-agent gaps, and different sensitivity to the stabilization mechanisms
introduced in this work. Extending the evaluation to more domains and model
scales is therefore an important direction for future work.
